\chapter{Будущее: память на клеточных комплексах}
\label{ch:cellcomplex}

Это приложение — про наиболее амбициозное структурное направление. Оно серьёзнее
гиперграфов (приложение~\ref{ch:hypergraphs}) и относится к области R\&D, но обсудить
его стоит, потому что оно даёт \emph{единую} картину.

\section{Идея: всё — объекты одного рода, разных размерностей}

В обычном графе узлы и рёбра — сущности \emph{разной} природы. Идея клеточного
комплекса: сделать их сущностями \emph{одного} рода, но разной \emph{размерности}.
\emph{Клетки}: 0-клетки — узлы (воспоминания, концепты), 1-клетки — рёбра (связи),
2-клетки — выводы из нескольких посылок, и так далее. Все они — «клетки», просто
живущие на разных этажах.

Связывает этажи \emph{граничный оператор} $\partial$: у каждой $k$-клетки есть
\emph{граница} — набор $(k{-}1)$-клеток, на которых она «натянута». Граница 1-клетки
(ребра) — два её узла; граница 2-клетки (вывода) — цикл/цепочка из посылок.

\begin{intuition}[Вывод рождается, когда построена его граница]
Твоя формулировка точна: вывод \emph{рождается}, когда его граница построена на
посылках. В терминах комплекса это \emph{приклеивание} $k$-клетки по её границе:
как только в памяти есть подходящая цепочка посылок (граница), поверх неё можно
«натянуть» клетку-вывод. Узел-вывод — это заполненная область, чей контур —
посылки.
\end{intuition}

\section{Почему это решает возражение к симплексам}

Здесь ключевой технический факт. Симплициальные комплексы требуют downward closure
(все грани обязаны существовать) — и потому плохи для разреженных n-арных выводов
(приложение~\ref{ch:hypergraphs}). \emph{Клеточные} (CW) комплексы этого
\textbf{не} требуют: $k$-клетку можно приклеить к границе \emph{произвольным}
приклеивающим отображением, не заполняя все промежуточные симплексы. То есть
клеточные комплексы \emph{снимают} ровно то возражение, из-за которого симплексы не
подходят, — и при этом дают единую алгебраическую структуру. Это делает их более
принципиальным ответом на n-арность, чем и симплексы, и даже голые гиперрёбра.

\section{Ходж-лапласианы: пропагация между этажами}

Граничные операторы $\partial_k$ порождают \emph{Ходж-лапласианы} на каждой
размерности:
\[
L_k=\partial_{k+1}\partial_{k+1}^\top+\partial_k^\top\partial_k .
\]
Это обобщение графового лапласиана (случай $k=0$) на клетки высших размерностей. Оно
даёт \emph{принципиальную} диффузию не только между узлами, но и между рёбрами,
между выводами, а также \emph{между этажами} (вверх и вниз по $\partial$). Иначе
говоря, вся динамика уверенности и активации (главы~\ref{ch:propagation},
\ref{ch:topology}) получает единое спектральное описание, охватывающее и выводы как
полноправные объекты.

\section{Прямая связь с чтением по DTI}

Это не абстрактная экзотика: ровно эта математика лежит в статьях, которые уже есть в
DTI-волте проекта. Bodnar и др. «WL Go Cellular» и «WL Go Topological» — передача
сообщений на клеточных и симплициальных комплексах; sheaf-лапласианы (cellular sheaf)
— их обогащение. Таким образом «память на клеточных комплексах» — очередной мост
между двумя проектами автора: инструменты топологического и клеточного обучения из
магистерской работы напрямую применимы к структуре памяти Рин.

\section{Честная оценка: северная звезда, не ближний шаг}

\begin{principle}[Клеточные комплексы — направление, а не спринт]
Это самая красивая и самая объединяющая картина: один род объектов, разные
размерности; вывод — высшая клетка; пропагация — Ходж-лапласиан по этажам. Но
выигрыш над прагматичной factor-кодировкой (приложение~\ref{ch:hypergraphs}) пока
\emph{не доказан}, а цена высока: динамические клеточные комплексы с приклеиванием и
отклеиванием клеток на мутирующей памяти, эффективные Ходж-лапласианы в реальном
времени — всё это серьёзный R\&D. Разумный порядок: ближний горизонт — factor-узлы;
средний — первоклассные гиперрёбра при необходимости; дальний/R\&D — клеточные
комплексы ради единого алгебраического описания.
\end{principle}

\begin{workedbox}[Как это выглядело бы у Рин]
Узлы «Португалия», «коты», «универ» — 0-клетки. Вывод «купи корм по пути» — 2-клетка,
приклеенная к границе из этих трёх посылок (без требования, чтобы посылки были
связаны попарно). Позже, если посылка «есть коты» опровергается, граница 2-клетки
рушится — и вывод \emph{автоматически} теряет опору (ретракция как разрушение
границы). Ходж-лапласиан при этом единообразно описывает, как понижение уверенности
течёт с этажа посылок на этаж вывода. Красиво — и именно поэтому стоит держать это
как цель, к которой архитектура может дозреть.
\end{workedbox}
